\pdfvariable suppressoptionalinfo 512
\pdfvariable trailerid {[<C3712026090400000000000000000000><C3712026090400000000000000000000>]}
\documentclass[10pt]{article}
\usepackage[margin=0.72in]{geometry}
\usepackage{amsmath,amssymb,amsthm,booktabs,array,microtype,xurl}
\usepackage{unicode-math}
\setmainfont{Latin Modern Roman}
\setmathfont{Latin Modern Math}
\providecommand{\CRevisionRound}{2}
\newtheorem{theorem}{Theorem}
\newtheorem{proposition}{Proposition}
\newtheorem{lemma}{Lemma}
\newcommand{\HH}{\mathsf H}
\newcommand{\PP}{\mathcal P}
\newcommand{\CC}{\mathcal C}
\newcommand{\Spec}{\operatorname{Spec}}
\newcommand{\Scope}{\texttt{NO\_BAD\_EULER\_OR\_ROOT\_NUMBER}}
\title{Rational-Flux Anisotropic Harper Dynamics:\
Chambers Phase Collapse and the Algebraic Bloch-Band Atlas}
\author{HCS-C371 / HEN-O355}
\date{4 September 2026}

\begin{document}
\maketitle

\begin{abstract}
For every reduced rational magnetic flux and every positive hopping
anisotropy, we fix one magnetic-cell phase convention and derive the exact
Chambers characteristic identity.  The proof exposes why the vertical
transfer trace has only three Laurent modes and fixes both phase signs.
\ifnum\CRevisionRound>0
We then recover the full two-dimensional spectrum as a polynomial preimage,
give the exact algebraic band-contact criterion, prove Aubry duality, flux
reversal and parity, and isolate the forced central contact at every even
denominator.  The one- and two-site cells are computed separately.
\fi
\ifnum\CRevisionRound>1
An exact cyclotomic lane and 74,880 independently reconstructed Bloch fibers
audit the result.  Harper--Chambers theory is established background, and
the final Route-A verdict remains negative despite natural quantization.
\fi
\end{abstract}

\section{Magnetic owner and phase convention}
Let $(p,q)=1$, $q\ge3$, and $\lambda>0$.  On
$\ell^2(\mathbb Z^2)$ consider the bounded self-adjoint operator
\begin{align}
 (\HH_{p/q,\lambda}\Psi)_{m,n}
 ={}&\Psi_{m+1,n}+\Psi_{m-1,n}\notag\\
 &+\lambda e^{2\pi i pm/q}\Psi_{m,n+1}
 +\lambda e^{-2\pi i pm/q}\Psi_{m,n-1}.
 \label{eq:lattice}
\end{align}
Fourier transformation in the second coordinate and magnetic Bloch
reduction give the $q$-dimensional fiber
\begin{equation}
 (\HH(k_x,k_y)u)_m=u_{m+1}+u_{m-1}
 +2\lambda\cos\!\left(k_y+\frac{2\pi pm}{q}\right)u_m,
 \qquad u_{m+q}=e^{iqk_x}u_m .
 \label{eq:fiber}
\end{equation}
Thus $X=qk_x$ is the \emph{total} horizontal boundary phase.  This
distinction is essential: throughout,
\begin{equation}
 D_{p/q,\lambda}(E;k_x,k_y)
 :=\det\!\left(EI_q-\HH_{p/q,\lambda}(k_x,k_y)\right).
 \label{eq:D}
\end{equation}
This manuscript's first revision is the \textbf{Chambers phase-collapse
owner}.

\section{The Chambers polynomial}
Put $\zeta=e^{2\pi i/q}$, $y=e^{ik_y}$, and
\begin{equation}
 A_m(E,y)=
 \begin{pmatrix}
 E-\lambda(y\zeta^{pm}+y^{-1}\zeta^{-pm})&-1\\
 1&0
 \end{pmatrix},
 \qquad M_q=A_{q-1}\cdots A_0 .
 \label{eq:transfer}
\end{equation}

\begin{theorem}[Anisotropic Chambers identity]\label{thm:chambers}
There is a unique real monic polynomial $\PP_{p/q,\lambda}$ of degree $q$
such that
\begin{equation}
 D_{p/q,\lambda}(E;k_x,k_y)
 =\PP_{p/q,\lambda}(E)-2\cos(qk_x)
 -2\lambda^q\cos(qk_y)
 \label{eq:chambers}
\end{equation}
for every real $E,k_x,k_y$.
\end{theorem}

\begin{proof}
Since $\det A_m=1$, the Floquet multipliers of $M_q$ are
$e^{\pm iqk_x}$.  Both sides below are monic of degree $q$ and have the
same roots, hence
\begin{equation}
 D_{p/q,\lambda}=\operatorname{tr}M_q-2\cos(qk_x).
 \label{eq:traceD}
\end{equation}
Replacing $y$ by $\zeta^p y$ cyclically permutes the potential sequence in
\eqref{eq:transfer}.  The monodromy changes by cyclic conjugacy, so its
trace is invariant.  Because $p$ is invertible modulo $q$, a Laurent
polynomial of $y$ with degrees between $-q$ and $q$ and this invariance can
use only the modes $-q,0,q$.

The coefficient of $y^q$ has one contribution: select
$-\lambda y\zeta^{pm}$ in every upper-left entry.  It equals
\begin{equation}
 (-\lambda)^q\zeta^{p q(q-1)/2}=-\lambda^q .
 \label{eq:extreme}
\end{equation}
Indeed, odd $q$ makes $q-1$ even, while even $q$ forces $p$ odd.  Complex
conjugation gives the same coefficient for $y^{-q}$.  The remaining mode
is independent of both phases, real, and monic; call it
$\PP_{p/q,\lambda}(E)$.  Equations \eqref{eq:traceD} and
\eqref{eq:extreme} prove \eqref{eq:chambers}.  Monicity gives uniqueness.
\end{proof}

The theorem is convention-complete: using the boundary character $X$ in
place of $k_x$ rewrites the horizontal term as $-2\cos X$, not
$-2\cos(qX)$.  This is precisely the convention tested by the executable
fibers.

\ifnum\CRevisionRound>0
\section{Spectrum and algebraic band contacts}
This revision is the \textbf{spectrum, duality, and edge owner}.  Write
\begin{equation}
 \CC_{q,\lambda}=2(1+\lambda^q),\qquad
 B_{p/q,\lambda}(E)=\PP_{p/q,\lambda}(E)^2-\CC_{q,\lambda}^2.
 \label{eq:edgepoly}
\end{equation}

\begin{lemma}[Sourced cyclic-continuant evaluation]\label{lem:central}
If $q$ is even, then
\begin{equation}
 \PP_{p/q,\lambda}(0)=2(-1)^{q/2}(1+\lambda^q).
 \label{eq:centralvalue}
\end{equation}
\end{lemma}

\begin{proof}
Lamoureux and Mingo use
$h_{\theta,L}=u+u^{-1}+(L/2)(v+v^{-1})$.  Their Theorem~2.5 proves the
vanishing of every intermediate cyclic-matching sum, and their
Corollary~2.6 gives, for even $q$,
\[
 \Delta_{p/q,L}(0)=2(-1)^{q/2}\bigl(1+(L/2)^q\bigr)
 \quad\cite{LamoureuxMingo}.
\]
Our diagonal potential is $2\lambda\cos(\cdot)$, so their parameter is
$L=2\lambda$.  Comparing the phase-independent monic term in their
Chambers formula with Theorem~\ref{thm:chambers} identifies
$\Delta_{p/q,2\lambda}=\PP_{p/q,\lambda}$: their defining vertical phase
$\pi/(2q)$ kills the explicit $q$th cosine, while our transfer trace then
equals $\PP$.  This also fixes our
$\det(EI-\HH)$ convention.  Substitution gives \eqref{eq:centralvalue}.
Thus this all-$q$ step is sourced rather than inferred from the finite
cyclotomic audit.
\end{proof}

\begin{theorem}[Complete Bloch spectral atlas]\label{thm:spectrum}
The full two-dimensional spectrum of \eqref{eq:lattice} is
\begin{equation}
 \Spec(\HH_{p/q,\lambda})
 =\{E\in\mathbb R:|\PP_{p/q,\lambda}(E)|\le\CC_{q,\lambda}\}.
 \label{eq:preimage}
\end{equation}
The two edge factors are actual fiber characteristic polynomials:
\begin{equation}
 \PP(E)-\CC_{q,\lambda}=D(E;0,0),\qquad
 \PP(E)+\CC_{q,\lambda}=D(E;\pi/q,\pi/q).
 \label{eq:edgefibers}
\end{equation}
Consequently the algebraic band-edge multiset is the real zero multiset of
$B$.  An
edge label is multiple precisely when
\begin{equation}
 \PP(E)=\pm\CC_{q,\lambda},\qquad \PP'(E)=0.
 \label{eq:contact}
\end{equation}
Moreover,
\begin{align}
 \PP_{p/q,\lambda}(E)
 &=\lambda^q\PP_{p/q,1/\lambda}(E/\lambda),
 \label{eq:duality}\\
 \PP_{p/q,\lambda}(E)&=\PP_{(q-p)/q,\lambda}(E),
 \label{eq:reversal}\\
 \PP_{p/q,\lambda}(-E)&=(-1)^q\PP_{p/q,\lambda}(E).
 \label{eq:parity}
\end{align}
If $q$ is even, then
\begin{equation}
 \PP_{p/q,\lambda}(0)=2(-1)^{q/2}(1+\lambda^q),
 \qquad \PP'_{p/q,\lambda}(0)=0,
 \label{eq:central}
\end{equation}
so zero is a multiple central edge.
\end{theorem}

\begin{proof}
Magnetic Bloch theory is a direct-integral decomposition into
\eqref{eq:fiber}.  By \eqref{eq:chambers}, $E$ lies in some fiber spectrum
exactly when
\[
 \PP(E)=2\cos(qk_x)+2\lambda^q\cos(qk_y).
\]
The phases vary independently, and their sum fills
$[-\CC_{q,\lambda},\CC_{q,\lambda}]$, proving \eqref{eq:preimage}.
Putting $(k_x,k_y)=(0,0)$ and $(\pi/q,\pi/q)$ in
\eqref{eq:chambers} proves \eqref{eq:edgefibers}.  Both endpoint fibers
are real symmetric: their total horizontal boundary phases are $0$ and
$\pi$, and their diagonal potentials are real.  Thus both edge factors
have only real roots, with multiplicity, and
$B=(\PP-\CC)(\PP+\CC)$ is real-rooted.  Equation \eqref{eq:contact} follows from
$B'=2\PP\PP'$ and $\PP\ne0$ at an edge.  It is an algebraic
coalescence criterion, not a claim that every remaining gap is open.

Exchange of the two lattice axes reverses the magnetic orientation and
swaps the hopping amplitudes.  After division by $\lambda$, magnetic
Fourier duality identifies the result with
$\lambda\HH_{(q-p)/q,1/\lambda}$.  Complex conjugation identifies the
fluxes $p/q$ and $(q-p)/q$.  Comparing the phase-independent monic terms
in \eqref{eq:chambers} proves \eqref{eq:duality}--\eqref{eq:reversal}.

The bipartite involution
$\Gamma\Psi_{m,n}=(-1)^{m+n}\Psi_{m,n}$ conjugates $\HH$ to $-\HH$ and
shifts both quasimomenta by $\pi$.  Substitution in
\eqref{eq:chambers} proves \eqref{eq:parity}.  For even $q$,
Lemma~\ref{lem:central} gives the first formula in \eqref{eq:central}.
Equation \eqref{eq:parity} makes $\PP$ even, so $\PP'(0)=0$.  This proves
the derivative formula without extrapolating the finite exact lane.
\end{proof}

\begin{proposition}[Accumulated small-cell boundaries]\label{prop:small}
When coincident wrap neighbors are added rather than overwritten,
Theorem~\ref{thm:chambers} has the direct boundary polynomials
\begin{equation}
 q=1,\ p=0:\quad \PP(E)=E,
 \qquad
 q=2,\ p=1:\quad \PP(E)=E^2-2(1+\lambda^2).
 \label{eq:small}
\end{equation}
The two $q=2$ folded bands meet at zero.  As $\lambda\downarrow0$, the
spectrum tends to the decoupled horizontal-chain interval $[-2,2]$;
reciprocal duality is not evaluated at $\lambda=0$.
\end{proposition}

\begin{proof}
For $q=1$ the scalar fiber is
$2\cos k_x+2\lambda\cos k_y$.  For $q=2$, write
$X=e^{2ik_x}$ and $Y=e^{ik_y}$.  The off-diagonal entries are
$1+X^{-1}$ and $1+X$, while the diagonal entries are
$\pm\lambda(Y+Y^{-1})$.  Its determinant is
\[
 E^2-2(1+\lambda^2)-(X+X^{-1})
 -\lambda^2(Y^2+Y^{-2}),
\]
which proves \eqref{eq:small} and the asserted contact.
\end{proof}
\fi

\ifnum\CRevisionRound>1
\section{Independent evidence}
This revision is the \textbf{evidence, collision, and route owner}.
The evidence exhausts all reduced fractions with $3\le q\le16$ and the
five anisotropies $1/2,2/3,1,3/2,2$.
\begin{table}[h]
\centering
\small
\begin{tabular}{@{}lr@{}}
\toprule
receipt & count\\
\midrule
reduced fluxes / parameter panels & $78$ / $390$\\
Bloch total-phase grid per panel & $12\times16$\\
Hermitian fibers / eigenvalues & $74{,}880$ / $825{,}600$\\
characteristic-determinant probes & $224{,}640$\\
exact cyclotomic reduced fluxes, $q\le10$ & $30$\\
\bottomrule
\end{tabular}
\caption{Finite regression scale.  Sampling is not the proof.}
\label{tab:evidence}
\end{table}

The producer obtains $\PP$ from transfer-polynomial multiplication.  The
checker never imports it: it reconstructs $\PP$ from the characteristic
polynomial of one Hermitian reference fiber and then rebuilds every matrix
in Table~\ref{tab:evidence}.  A third lane works exactly in
$\mathbb Q[\zeta_q]/(\Phi_q)$ with a symbolic anisotropy degree; it checks
the Laurent support, both extreme coefficients, duality, reversal, parity,
and \eqref{eq:central} through $q=10$.  It is a regression check, not the
all-denominator proof supplied by Lemma~\ref{lem:central}.  Isolated replay,
repaired-hash mutations and strict JSON/YAML parsers close the computational
boundary.

\section{Source and collision audit}
Harper's magnetic difference equation, Chambers' rational-flux relation,
and Hofstadter's magnetic bands are classical
\cite{Harper,Chambers,Hofstadter,Avron,JKK}.  Lamoureux and Mingo directly
precede the cyclic-matching cancellation and even-denominator constant-term
formula used in Lemma~\ref{lem:central}~\cite{LamoureuxMingo}.  We claim no
priority for these objects or formulas.  The package contribution is a
convention-locked reconstruction, the explicit real endpoint-fiber edge
factorization \eqref{eq:edgefibers}, and an independently replayable atlas.

The direct workspace neighbor is HCS-C15/HEN-O30.  It owns a critical
Weyl--Harper block $U+U^*+V+V^*$ at flux $1/3^m$ and the return of its top
spectral edge along a Heisenberg congruence tower.  It does not own the
two-phase identity \eqref{eq:chambers}, all reduced fluxes, anisotropic
duality, or the edge polynomial \eqref{eq:edgepoly}.  C293 is a magnetic
Grushin cylinder, C340 is a Lam\'e Floquet operator, and C356 is a QWZ Chern
pump.  No Chern, transport, edge-state, irrational-flux Cantor-spectrum, or
Ten-Martini result is imported here.

\section{Route decision and nonclaims}
Reduced rational flux gives exact cyclotomic magnetic translations, earning
$A0_{\rm WEAK\_ARITHMETIC\_RELATION}$.  It does not give rational-prime
ownership, prime-power repetition, an isolated primitive-orbit ledger, or
a logarithmic prime clock, so $A1_{\rm FAIL}$.  The finite source
characteristic polynomial is neither a target Euler product nor a target
Fredholm determinant, so $A2_{\rm FAIL}$; no target continuation,
functional equation, divisor or zero-counting law exists, so
$A3_{\rm FAIL}$.  Equation \eqref{eq:lattice} is already a bounded
self-adjoint Hamiltonian, hence $A4_{\rm NATURAL\_QUANTIZATION}$, but its
Bloch bands are not target zeros or a Hilbert--P\'olya realization.

The overall verdict is \texttt{ROUTE\_A\_REJECTED} under the literal scope
\Scope.  No target arithmetic local datum, Euler factor, root number,
automorphy, target-zero match, or Hilbert--P\'olya operator is asserted.
Route B remains locked.
\fi

\begin{thebibliography}{9}\small
\bibitem{Harper} P. G. Harper,
\emph{Single band motion of conduction electrons in a uniform magnetic
field}, Proc. Phys. Soc. A 68 (1955), 874--878.
\bibitem{Chambers} W. G. Chambers,
\emph{Linear-network model for magnetic breakdown in two dimensions},
Phys. Rev. 140 (1965), A135--A143.
\bibitem{Hofstadter} D. R. Hofstadter,
\emph{Energy levels and wave functions of Bloch electrons in rational and
irrational magnetic fields}, Phys. Rev. B 14 (1976), 2239--2249.
\bibitem{Avron} J. Avron, P. H. M. van Mouche, and B. Simon,
\emph{On the measure of the spectrum for the almost Mathieu operator},
Comm. Math. Phys. 132 (1990), 103--118.
\bibitem{LamoureuxMingo} M.~P. Lamoureux and J.~A. Mingo,
\emph{On the characteristic polynomial of the almost Mathieu operator},
Proc. Amer. Math. Soc. 135 (2007), 3205--3215.
\url{https://doi.org/10.1090/S0002-9939-07-08830-2}.
\bibitem{JKK} S. Jitomirskaya, L. Konstantinov, and I. Krasovsky,
\emph{On the spectrum of critical almost Mathieu operators in the rational
case}, arXiv:2007.01005.
\end{thebibliography}
\end{document}
